\section{Proof of the Reduction-Induced Interaction Theorem}
\label{app:interaction-proof}

For self-containment, we reproduce the P1 proof of
\cref{thm:reduction-induced-interaction} and
\cref{prop:reduction-mobius-noncommutation}, then prove the new observable
readout law in \cref{thm:readout-interaction-transfer}.

\subsection{Continuity and smoothness of the minimizer section}

Fix $z_0=(y_0,u_0)$ and write $\sigma_0=\sigma_\star(z_0)$.  Let $z_n\to z_0$.  Evaluating each fiber objective at $\sigma_0$ shows that $E(z_n,\sigma_\star(z_n))$ is bounded above for large $n$.  Local uniform inf-compactness therefore places the minimizers in one relatively compact set.  Every subsequence has a further convergent subsequence, say $\sigma_\star(z_{n_j})\to\bar\sigma$.

For any fixed $\sigma$, optimality gives
\[
E(z_{n_j},\sigma_\star(z_{n_j}))\le E(z_{n_j},\sigma).
\]
Continuity and passage to the limit yield $E(z_0,\bar\sigma)\le E(z_0,\sigma)$ for all $\sigma$.  Uniqueness at $z_0$ implies $\bar\sigma=\sigma_0$.  Every convergent subsequence has this limit, hence the full minimizer sequence converges.  Thus $\sigma_\star$ is continuous.

The interior first-order condition is
\begin{equation}
E_\sigma(z,\sigma_\star(z))=0.
\label{eq:appendix-stationarity}
\end{equation}
At every minimizer, $D_\sigma E_\sigma=H$ is invertible.  The implicit-function theorem gives a local $C^r$ critical section through each $(z,\sigma_\star(z))$.  Continuity of the global minimizer section makes it coincide locally with that critical section.  These local representations agree by uniqueness, so $\sigma_\star$ is globally $C^r$ on the parameter domain.

\subsection{Hidden response and effective Hessian}

Differentiate \cref{eq:appendix-stationarity} with respect to $z=(y,u)$:
\[
E_{\sigma z}+H D_z\sigma_\star=0.
\]
Since $H$ is invertible,
\[
D_z\sigma_\star=-H^{-1}E_{\sigma z},
\]
which proves \cref{eq:hidden-response-law}.

For the reduced value, the chain rule and stationarity give the envelope identity
\[
D_z\bar E
=E_z+E_\sigma D_z\sigma_\star
=E_z.
\]
Differentiating again,
\[
D^2_{zz}\bar E
=E_{zz}+E_{z\sigma}D_z\sigma_\star
=E_{zz}-E_{z\sigma}H^{-1}E_{\sigma z}.
\]
This is \cref{eq:effective-schur-hessian}.  At a critical hidden point the displayed quadratic form is the Schur complement of the vertical Hessian, hence is invariant under smooth changes of hidden coordinates.

\subsection{Affine interventions and interaction curvature}

Under \cref{eq:affine-interventions}, $E_{uu}=0$ and $E_{\sigma u}=G$.  The $u$ block of \cref{eq:effective-schur-hessian} is therefore
\[
D^2_{uu}\bar E=-G^*H^{-1}G.
\]
For every $a\in\R^m$,
\[
a^\top D^2_{uu}\bar E\,a
=-\ip{Ga}{H^{-1}Ga}\le0,
\]
because $H^{-1}$ is positive definite.  This proves \cref{eq:interaction-curvature}.  Componentwise,
\begin{equation}
\partial^2_{u_i u_j}\bar E
=-\ip{D_\sigma\Delta_i}{H^{-1}D_\sigma\Delta_j}.
\label{eq:appendix-component-curvature}
\end{equation}

The same concavity conclusion holds without differentiability whenever the infimum is finite: for fixed $(y,\sigma)$, the unreduced energy is affine in $u$, and a pointwise infimum of affine functions is concave.

\subsection{Observable readout interaction}

Fix $y$ and suppress it from the notation.  Write
$v_i=\partial_{u_i}\sigma_\star$.  Differentiating stationarity once gives
\[
Hv_i+E_{\sigma u_i}=0,
\qquad
v_i=-H^{-1}E_{\sigma u_i}.
\]
Differentiate this identity with respect to $u_j$ along the minimizing
section.  The total derivative of $H=E_{\sigma\sigma}$ contributes
$E_{\sigma\sigma\sigma}[v_j,v_i]+E_{\sigma\sigma u_j}v_i$, while the total
derivative of $E_{\sigma u_i}$ contributes
$E_{\sigma\sigma u_i}v_j+E_{\sigma u_i u_j}$.  Therefore
\[
H\partial^2_{u_i u_j}\sigma_\star
+E_{\sigma\sigma\sigma}[v_i,v_j]
+E_{\sigma\sigma u_i}v_j
+E_{\sigma\sigma u_j}v_i
+E_{\sigma u_i u_j}=0.
\]
Invertibility of $H$ proves \cref{eq:second-hidden-response}.

For $F_R(u)=R(u,\sigma_\star(u))$, the first derivative is
\[
\partial_{u_i}F_R=R_{u_i}+R_\sigma v_i.
\]
Differentiating with respect to $u_j$ and applying the product and chain rules
gives
\[
\partial^2_{u_i u_j}F_R
=R_{u_i u_j}+R_{u_i\sigma}v_j+R_{\sigma u_j}v_i
+R_{\sigma\sigma}[v_i,v_j]
+R_\sigma\partial^2_{u_i u_j}\sigma_\star,
\]
which is \cref{eq:readout-interaction-transfer}.  Applying the two coordinate
finite-difference operators and the fundamental theorem of calculus proves
\cref{eq:readout-mobius-integral}.  If $E$ is affine in $u$, direct
differentiation gives $E_{\sigma u_i u_j}=0$ and
$E_{\sigma\sigma u_i}=D^2_{\sigma\sigma}\Delta_i$.  No sign follows for the
readout Hessian because its five terms need not form a negative Gram matrix.

\subsection{M\"obius effects as integrated derivatives}

For a coordinate $i$, let $\mathsf T_i$ replace $u_i=0$ by $u_i=1$.  The fundamental theorem of calculus gives
\[
(\mathsf T_i-I)\bar E(0)
=\int_0^1\partial_{u_i}\bar E(t_ie_i)\,\dd t_i.
\]
The coordinate difference operators commute.  Repeating the identity for every $i\in T$ gives
\[
\prod_{i\in T}(\mathsf T_i-I)\bar E(0)
=\int_{[0,1]^T}
\partial_T^{|T|}\bar E\left(\sum_{i\in T}t_ie_i\right)
\prod_{i\in T}\dd t_i.
\]
Expanding the product on the left yields
\[
\sum_{B\subseteq T}(-1)^{|T|-|B|}\bar E(\mathbf1_B)=\zeta_T,
\]
proving \cref{eq:mobius-derivative-integral}.  For $T=\{i,j\}$, substitution of \cref{eq:appendix-component-curvature} proves \cref{eq:mobius-schur-bridge}.

A zero integral does not force the integrand to vanish because mixed curvature may change sign.  Pointwise vanishing is equivalent to inverse-Hessian orthogonality of $D_\sigma\Delta_i$ and $D_\sigma\Delta_j$.

\subsection{Noncommutation and analytical examples}

For \cref{prop:reduction-mobius-noncommutation}, let
\[
E_0(\sigma)=\sigma^2,
\qquad
E_1(\sigma)=(\sigma-1)^2.
\]
Both reduced values are zero, so their reduced first effect is zero.  The pointwise first difference is $E_1-E_0=1-2\sigma$, whose infimum is $-\infty$.  Thus no general exchange law exists.

If instead $E_A(\sigma)=E_0(\sigma)+\sum_{i\in A}c_i$ with fiber constants $c_i$, then the constants leave the minimizer unchanged and move outside the infimum.  The reduced response is additive, with singleton effects $c_i$ and all higher effects zero.

Finally, for
\[
E(\sigma,u)=\frac h2\sigma^2+u_1a\sigma+u_2b\sigma,
\]
stationarity yields $\sigma_\star=-(au_1+bu_2)/h$.  Substitution gives
\[
\bar E(u)=-\frac{(au_1+bu_2)^2}{2h}.
\]
The Boolean pair contrast is
\[
\bar E(1,1)-\bar E(1,0)-\bar E(0,1)+\bar E(0,0)
=-\frac{ab}{h},
\]
matching \cref{eq:mobius-schur-bridge}.

For the coupled log-preconditioner in
\cref{eq:coupled-log-preconditioner}, $h>|\rho|$ makes
\[
H=\begin{pmatrix}h&\rho\\ \rho&h\end{pmatrix}
\quad\text{positive definite},\qquad
H^{-1}=\frac1{h^2-\rho^2}
\begin{pmatrix}h&-\rho\\-\rho&h\end{pmatrix}.
\]
Stationarity gives
\[
H\sigma-b(g)+u=0,
\qquad
\sigma_\star(g,u)=H^{-1}(b(g)-u).
\]
Completing the square yields the reduced scalar calibration response
\[
\bar E(g,u)
=-\frac12(b(g)-u)^\top H^{-1}(b(g)-u).
\]
Because $D_\sigma\Delta_1=e_1$, $D_\sigma\Delta_2=e_2$, and
$G=I_2$, its mixed derivative is constant:
\[
\partial_{u_1u_2}^2\bar E
=-e_1^\top H^{-1}e_2
=\frac{\rho}{h^2-\rho^2}.
\]
The integral over the Boolean intervention square has the same value, proving
\cref{eq:coupled-log-preconditioner-pair}.  The selected update
$q(\sigma_\star,g)=-\diag(e^{-\sigma_{\star,1}/2},
e^{-\sigma_{\star,2}/2})g$ is a positive diagonal preconditioned gradient
step for every intervention amplitude.

Let $D=h^2-\rho^2$ and $\sigma_\star(0)=H^{-1}b(g)$.  The first hidden
coordinate is
\[
\sigma_{\star,1}(u)
=\sigma_{\star,1}(0)-\frac hD u_1+\frac\rho D u_2.
\]
Substitution into $q_1=-g_1e^{-\sigma_{\star,1}/2}$ yields the exponential
readout displayed in the main text.  For a function
$q_1(0)e^{a u_1+c u_2}$, the Boolean pair difference is exactly
$q_1(0)(e^a-1)(e^c-1)$.  Taking $a=h/(2D)$ and
$c=-\rho/(2D)$ proves \cref{eq:coupled-update-readout-pair} and completes the
optimizer realization.
